\section*{Chapitre \ref{ch:b3:endocrinology} --- Endocrinologie et homéostasie}

\begin{solution}{exo:b3:endocrinology:1}
\omterm{def:b3:endocrinology:glucose}{Insuline}~: peptide, récepteur de surface (une tyrosine kinase), agit en
quelques minutes, demi-vie d'environ \qty{5}{min}. Cortisol~: stéroïde,
\omterm{def:b3:endocrinology:hormone}{récepteur nucléaire}, quelques heures, demi-vie \qty{80}{min} (surtout
lié à un porteur). Adrénaline~: amine, récepteurs de surface couplés aux
protéines~G, quelques secondes, demi-vie d'environ \qty{2}{min}.
Thyroxine~: acide aminé iodé, \omterm{def:b3:endocrinology:hormone}{récepteur nucléaire}, quelques jours,
demi-vie \qty{7}{d} (liée à des porteurs). Vasopressine~: peptide,
récepteur de surface couplé aux protéines~G (AMPc dans le \omterm{def:b3:renal-osmoregulation:nephron}{tube
collecteur}), quelques minutes, demi-vie d'environ \qty{15}{min}.
\end{solution}

\begin{solution}{exo:b3:endocrinology:2}
\omterm{def:b3:endocrinology:axes}{Hypothalamus} (TRH) $\to$ \omterm{def:b3:endocrinology:axes}{hypophyse} (TSH) $\to$ thyroïde (T$_{4}$), la
T$_{4}$ inhibant les deux étages du dessus. Thyroïde détruite~: T$_{4}$
basse, TSH élevée. \omterm{def:b3:cancer-biology:hallmarks}{Tumeur} sécrétant la TSH~: TSH élevée et T$_{4}$
élevée (le rétrocontrôle n'étouffe pas la \omterm{def:b3:cancer-biology:hallmarks}{tumeur}). Excès de comprimés de
thyroxine~: T$_{4}$ élevée, TSH étouffée. Carence en iode~: T$_{4}$
basse ou normale basse, TSH élevée, et la glande grossit sous elle en
goitre.
\end{solution}

\begin{solution}{exo:b3:endocrinology:3}
Le cerveau meurt d'une glycémie basse en quelques minutes et n'est
abîmé par une glycémie haute qu'au fil des ans~; et l'évolution a
rencontré la famine bien plus souvent que les festins. La défense contre
une baisse est donc redondante (\omterm{def:b3:endocrinology:glucose}{glucagon}, adrénaline, cortisol, \omterm{def:b3:endocrinology:hormone}{hormone}
de croissance) et la défense contre une hausse est unique. Prédiction~:
trop d'\omterm{def:b3:endocrinology:glucose}{insuline} est létal dans l'immédiat (le coma hypoglycémique), trop
peu est une maladie chronique --- ce que voit la clinique.
\end{solution}

\begin{solution}{exo:b3:endocrinology:4}
Un \omterm{def:b3:endocrinology:clock}{zeitgeber} est un repère de l'environnement qui entraîne une horloge~:
la lumière, l'heure des repas, la température, l'exercice, le rythme
social. L'horloge maîtresse emploie la lumière, par les \omterm{def:b3:sensory-systems:retina}{cellules
ganglionnaires} à mélanopsine de la \omterm{def:b3:sensory-systems:retina}{rétine}. Après huit fuseaux horaires
l'horloge ne se décale que d'environ une heure par jour, de sorte que
pendant une semaine le sommeil, le cortisol, la température et l'appétit
tournent à l'ancienne heure tandis que les horloges périphériques,
remises à l'heure par les repas, dérivent à leur propre rythme~: le
décalage horaire, pire vers l'est, puisque avancer l'horloge est plus
difficile que la retarder.
\end{solution}

\begin{solution}{exo:b3:endocrinology:5}
Réponse à mi-hauteur pour $250$ récepteurs occupés~: $H = K_{d}\times
250/(\num{20000} - 250) = \qty{1}{nM}\times 0.0127 = \qty{12.7}{pM}$,
quatre-vingts fois au-dessous de $K_{d}$. Avec \num{2000} récepteurs~: $250/1750
\times \qty{1}{nM} = \qty{143}{pM}$, onze fois moins sensible. Une
\omterm{def:b3:endocrinology:glucose}{insuline} durablement élevée régule négativement ses récepteurs, de sorte
qu'un effet donné demande plus d'\omterm{def:b3:endocrinology:glucose}{insuline}~: une composante de la
résistance à l'\omterm{def:b3:endocrinology:glucose}{insuline}, qui se nourrit d'elle-même.
\end{solution}

\begin{solution}{exo:b3:endocrinology:6}
$T_{4} = 13$~: $1.5\,\mathrm{e}^{0.92} = \qty{3.8}{mU/L}$~; $11$~:
$1.5\,\mathrm{e}^{1.84} = \qty{9.4}{mU/L}$~; $9$~: $1.5\,\mathrm{e}^{2.76}
= \qty{24}{mU/L}$~; $25$~: $1.5\,\mathrm{e}^{-4.6} = \qty{0.015}{mU/L}$. À
\qty{11}{pmol/L} la T$_{4}$ est encore dans \qtyrange{10}{20}{} alors
que la TSH fait plus du double de sa limite haute (hypothyroïdie
infraclinique)~; à $13$ la TSH est tout juste dans la plage~; à $9$ les
deux sont anormales~; à $25$ les deux le sont dans l'autre sens.
\end{solution}

\begin{solution}{exo:b3:endocrinology:7}
$L = s\beta/(a\gamma) = 0.2\times 0.05/(0.02\times 0.1) = 5$. Excès
stationnaire $g^{*} = (u/a)/(1 + L) = (2/0.02)/6 = \qty{16.7}{mg/dL}$~;
sans rétrocontrôle $u/a = \qty{100}{mg/dL}$. Excès d'\omterm{def:b3:endocrinology:glucose}{insuline} $i^{*} = \beta g^{*}/
\gamma = 0.05\times 16.7/0.1 = \qty{8.3}{mU/L}$.
\end{solution}

\begin{solution}{exo:b3:endocrinology:8}
Discriminant $(a - \gamma)^{2} - 4s\beta = 0.0064 - 0.04 < 0$~:
oscillatoire. $\omega = \sqrt{a\gamma + s\beta - (a+\gamma)^{2}/4} =
\sqrt{0.002 + 0.01 - 0.0036} = \qty{0.092}{min^{-1}}$, période $2\pi/
\omega = \qty{68}{min}$~; l'amplitude baisse d'un facteur $\mathrm{e}$ en $2/(a +
\gamma) = \qty{16.7}{min}$. En divisant $s$ par deux~: $s\beta = 0.005$,
discriminant $0.0064 - 0.02 < 0$, toujours oscillatoire, $\omega = \sqrt{0.007 - 0.0036}
= \qty{0.058}{min^{-1}}$, période \qty{108}{min}~; le temps de
décroissance est inchangé~; $L = 2.5$. Une boucle plus faible revient
plus lentement et tient une erreur stationnaire plus grande.
\end{solution}

\begin{solution}{exo:b3:endocrinology:9}
Pendant un an la prednisolone a étouffé la CRH et l'ACTH~; le cortex
surrénalien, non stimulé, s'est atrophié. À l'arrêt il n'y a plus de
stéroïde exogène, l'\omterm{def:b3:endocrinology:axes}{hypothalamus} et l'\omterm{def:b3:endocrinology:axes}{hypophyse} mettent des semaines à
repartir, et la surrénale rétrécie ne saurait répondre à l'ACTH si elle
venait. Une infection réclame dix fois le cortisol normal~; sans lui,
les vaisseaux se dilatent, la pression et le glucose tombent~: c'est
l'insuffisance surrénalienne aiguë. Mesuré~: ACTH basse (la lésion est
centrale), cortisol bas, et une faible montée du cortisol à l'ACTH
injectée (la glande a fondu) --- contrairement à la maladie d'Addison,
où l'ACTH est élevée. D'où la décroissance lente des doses.
\end{solution}

\begin{solution}{exo:b3:endocrinology:10}
À une variable~: $\dot x = f(x)$ avec $f' < 0$ a un point fixe unique
$x^{*}$~; pour $x > x^{*}$, $\dot x < 0$ et $x$ décroît de façon
monotone vers $x^{*}$, sans jamais le traverser (unicité des solutions),
et symétriquement en dessous~: pas d'oscillation. À deux variables~: la
matrice peut avoir des valeurs propres complexes, comme dans le
théorème, quand l'effecteur agit par une seconde variable ayant sa
propre échelle de temps. Une répression instantanée se poserait à un
niveau stationnaire~; le retard entre la transcription et la répression
nucléaire (traduction, dimérisation, phosphorylation, import) laisse la
protéine dépasser avant que le gène ne soit éteint et passer en dessous
avant qu'il ne soit rallumé, et la période vaut à peu près le double de
la somme du retard et du temps de dégradation de la protéine. Une
dégradation plus rapide de PER raccourcit la période~: une mutation
déstabilisante de PER2 humaine donne une horloge d'environ 20 heures et
une famille qui s'endort à 19 heures.
\end{solution}

\begin{solution}{exo:b3:endocrinology:11}
Exposant $k(\text{Ca} - 1.2)$~: à $1.10$, $\mathrm{e}^{-2}$, PTH $= 1/
(1 + 0.135) = 0.88$~; $1.15$~: $0.73$~; $1.20$~: $0.50$~; $1.25$~: $1/(1 +
2.72) = 0.27$~; $1.30$~: $1/(1 + 7.39) = 0.12$. Pente au \omterm{thm:b3:endocrinology:loop}{point de consigne}
$-k/4 = -5$ par mmol/L~: \qty{5}{\%} du maximum pour chaque
\qty{0.01}{mmol/L}. Le calcium doit être tenu à quelques pour cent près
et son effecteur agit sur un énorme tampon (l'os) sans dépassement, de
sorte qu'une réponse raide est sûre et nécessaire. Une boucle du glucose
aussi raide, agissant par une \omterm{def:b3:endocrinology:glucose}{insuline} ayant son propre temps
d'élimination, aurait un grand \omterm{thm:b3:endocrinology:loop}{gain de boucle} et oscillerait profondément
en hypoglycémie après chaque repas.
\end{solution}

\begin{solution}{exo:b3:endocrinology:12}
La glycation est irréversible et lente, de sorte qu'un globule d'âge
$\tau$ en a une fraction $kG\tau$ de glyquée~; moyennée sur des globules
d'âges uniformément répartis de \numrange{0}{120} jours, l'HbA1c
$= kG\times\qty{60}{d}$, proportionnelle à la glycémie moyenne.
Étalonnage~: $5.5\times 60k = 0.05$ donne $k =
\qty{1.5e-4}{}$ par (mmol/L)(jour)~; à \qty{10}{mmol/L}~: \qty{9.1}{\%}.
Avec des globules vivant 60 jours l'âge moyen est de 30 jours et l'HbA1c
est divisée par deux pour la même glycémie~: un patient à
\qty{10}{mmol/L} lirait \qty{4.5}{\%}, une valeur normale.
\end{solution}

\begin{solution}{pb:b3:endocrinology:1}
\textbf{1.} \qty{90}{mg/dL} $= \qty{0.9}{g/L} = 0.9/180 = \qty{5.0}{mmol/L}$~;
$0.9\times 15 = \qty{13.5}{g}$ dans le volume de distribution.
\textbf{2.} $75/15 = \qty{5}{g/L}$~: une hausse de \qty{500}{mg/dL},
\qty{28}{mmol/L}. Le pic réel est bien plus bas parce que l'absorption
s'étale sur une heure tandis que l'élimination retire le glucose à
mesure qu'il arrive, et que le foie en prend une large part au premier
passage.
\textbf{3.} $\qty{75}{g}/\qty{60}{min} = \qty{1.25}{g/min}$~; sur
\qty{15}{L}~: \qty{8.3}{mg/dL/min}, quatre fois $u_{0}$.
\textbf{4.} Excès stationnaire $u/a = 8.3/0.02 = \qty{417}{mg/dL}$,
constante de temps $1/a = \qty{50}{min}$~; après une heure $417(1 -
\mathrm{e}^{-1.2}) = \qty{291}{mg/dL}$ au-dessus du jeûne~: environ
\qty{380}{mg/dL}, \qty{21}{mmol/L}.
\textbf{5.} $L = 0.2\times 0.05/(0.02\times 0.1) = 5$~; excès stationnaire
$417/6 = \qty{69}{mg/dL}$, six fois moindre que sans \omterm{def:b3:endocrinology:glucose}{insuline}.
\textbf{6.} Le cerveau brûle \qty{120}{g} de glucose par jour, n'en
stocke pas et ne peut guère employer autre chose~: une glycémie basse
donne confusion et coma en quelques minutes. Une glycémie haute glyque
les protéines et, au fil des ans, abîme les petits vaisseaux de la
\omterm{def:b3:sensory-systems:retina}{rétine}, du rein et des nerfs, ainsi que les grosses artères.
\textbf{7.} $\lambda^{2} + (a + \gamma)\lambda + (a\gamma + s\beta) = 0$~:
trace $-0.12$, déterminant $0.002 + 0.01 = 0.012$, discriminant
$0.0144 - 0.048 = -0.0336$.
\textbf{8.} $\omega = \sqrt{0.0336}/2 = \qty{0.092}{min^{-1}}$~; période
\qty{68}{min}~; temps de décroissance $2/0.12 = \qty{16.7}{min}$.
\textbf{9.} $\dot g(0) = 100(-0.06 + c\,\omega) = -2$, donc $c\,\omega =
0.04$, $c = 0.44$.
\textbf{10.} $g = 0$ quand $\tan\omega t = -1/c = -2.29$, $\omega t =
\pi - 1.16 = 1.98$, $t = \qty{21.6}{min}$. Au premier minimum, vers
\qty{32}{min}~: $100\,\mathrm{e}^{-1.94}(\cos 2.97 + 0.44\sin 2.97) =
14.4\times(-0.91) \approx -\qty{13}{mg/dL}$, soit un glucose de \qty{77}{mg/dL}.
\textbf{11.} À $t = 0$, $\dot i = 100\beta > 0$~: l'\omterm{def:b3:endocrinology:glucose}{insuline} monte tant
que $g > 0$, et culmine quand $\beta g = \gamma i$, ce qui exige que $g$
soit encore positif --- donc le pic vient avant que le glucose atteigne
le jeûne, et après qu'il a commencé à baisser (il baisse dès $t = 0$).
Dans le modèle le pic est vers \qty{12}{min}.
\textbf{12.} Avec une entrée étalée sur une heure le pic survient quand
l'absorption et l'élimination s'équilibrent, plus bas et plus tard (30 à
60 min)~; le retour attend la fin de l'absorption, d'où les deux
heures~; le dépassement d'\omterm{def:b3:endocrinology:glucose}{insuline} est moindre pour une entrée
progressive, donc le creux est léger.
\textbf{13.} $L = 2.5$. Sain $g^{*} = (2/0.02)/6 = \qty{16.7}{mg/dL}$~;
résistant $100/3.5 = \qty{28.6}{mg/dL}$~: une hausse d'environ
\qty{12}{mg/dL} (\qty{0.7}{mmol/L}), glycémie à jeun voisine de
\qty{102}{mg/dL}, \qty{5.7}{mmol/L} --- le seuil de l'hyperglycémie
modérée à jeun.
\textbf{14.} $i^{*} = \beta g^{*}/\gamma$~: sain $8.3$, résistant
\qty{14.3}{mU/L}, rapport $1.7$~: une résistance à l'\omterm{def:b3:endocrinology:glucose}{insuline} compensée,
hyperinsulinémie avec une glycémie presque normale.
\textbf{15.} $L = 0.1\times 0.01/0.002 = 0.5$~; $g^{*} = 100/1.5 =
\qty{66.7}{mg/dL}$, \qty{50}{mg/dL} au-dessus de l'état sain~:
\qty{2.8}{mmol/L}, glycémie à jeun \qty{7.8}{mmol/L} (\qty{140}{mg/dL}),
au-dessus du seuil diabétique.
\textbf{16.} Discriminant $(a - \gamma)^{2} - 4s\beta = 0.0064 - 0.004 =
0.0024 > 0$~: retour monotone. $\lambda = -0.06 \pm 0.0245$~: $-0.0355$
et $-0.0845$ par minute~; constante de temps la plus lente \qty{28}{min},
contre \qty{16.7}{min} pour la décroissance saine~: la boucle
défaillante revient plus lentement et ne passe jamais dessous.
\textbf{17.} Glucose élevé avec \omterm{def:b3:endocrinology:glucose}{insuline} élevée~: résistance avec
compensation partielle --- un type~2~; la valeur à deux heures de
\qty{13}{mmol/L} dépasse \qty{11.1}{}, donc c'est un diabète, et non une
simple intolérance.
\textbf{18.} Pas de peptide~C signifie pas d'\omterm{def:b3:endocrinology:glucose}{insuline} endogène~: un
type~1. Sans \omterm{def:b3:endocrinology:glucose}{insuline} le muscle et le tissu adipeux ne peuvent capter le
glucose, la graisse est dégradée en acides gras et en corps cétoniques,
la protéine musculaire est catabolisée pour la néoglucogenèse, et le
glucose est perdu dans l'urine avec ses calories~: le patient meurt de
faim au milieu de l'abondance.
\textbf{19.} $5\times 9/5.5 = \qty{8.2}{\%}$.
\textbf{20.} $1.5\,\mathrm{e}^{0.46\times 3} = 1.5\times 3.97 =
\qty{6.0}{mU/L}$. T$_{4}$ dans la plage, TSH au-dessus~: hypothyroïdie
infraclinique, la glande défaille et l'\omterm{def:b3:endocrinology:axes}{hypophyse} compense.
\textbf{21.} $1.5\,\mathrm{e}^{0.46\times 8} = 1.5\times 39.6 =
\qty{59}{mU/L}$. Une T$_{4}$ basse avec une TSH de seulement $0.3$
signifie que l'\omterm{def:b3:endocrinology:axes}{hypophyse} ne répond pas~: la lésion est centrale
(\omterm{def:b3:endocrinology:axes}{hypophyse} ou \omterm{def:b3:endocrinology:axes}{hypothalamus}), non thyroïdienne.
\textbf{22.} Une demi-vie~: le niveau tombe à \qty{50}{\%}, et les
symptômes s'installent sur des semaines. Le cortisol, de demi-vie
\qty{80}{min}, a disparu en quelques heures après une dose manquée, et
un stress ce jour-là ne rencontre aucune défense~: la semaine de
thyroxine manquée est inconfortable, le jour de cortisol manqué peut
tuer.
\textbf{23.} T$_{4}$ élevée (l'\omterm{def:b3:adaptive-immunity:antibody}{anticorps} commande la glande sans égard
au rétrocontrôle)~; TSH étouffée par la T$_{4}$ élevée~; la relation
log-linéaire la rend indétectable --- à $T_{4} = 30$, $1.5\,\mathrm{e}^{-6.9}
= \qty{0.0015}{mU/L}$ --- de sorte qu'une TSH indétectable est le
premier signe.
\textbf{24.} Le niveau d'une \omterm{def:b3:endocrinology:hormone}{hormone} rapporte l'équilibre entre l'ordre
qu'elle reçoit et sa glande~: une TSH élevée avec une T$_{4}$ basse est
une glande défaillante et une TSH élevée avec une T$_{4}$ élevée une
\omterm{def:b3:endocrinology:axes}{hypophyse} emballée~; une \omterm{def:b3:endocrinology:glucose}{insuline} élevée avec un glucose élevé est un
corps résistant et une \omterm{def:b3:endocrinology:glucose}{insuline} basse avec un glucose élevé un îlot
détruit. Lue seule, une T$_{4}$ de \qty{7}{pmol/L} ou une \omterm{def:b3:endocrinology:glucose}{insuline} de
\qty{30}{mU/L} ne dit pas où est la faute~; lue avec celle qui la
commande, elle le dit.
\textbf{25.} \omterm{thm:b3:endocrinology:loop}{Gain de boucle} $L = 5$~; période \qty{68}{min} et temps de
décroissance \qty{17}{min}~; glycémie à jeun de la boucle défaillante \qty{7.8}{mmol/L}.
\end{solution}
